\documentclass[11pt]{article}
\usepackage[a4paper,margin=22mm]{geometry}
\usepackage{fontspec}
\setmainfont{DejaVu Serif}
\setsansfont{DejaVu Sans}
\usepackage{microtype}
\usepackage{xcolor}
\usepackage{hyperref}
\usepackage{enumitem}
\usepackage{parskip}
\definecolor{Accent}{HTML}{971D32}
\definecolor{Muted}{HTML}{666666}
\hypersetup{colorlinks=true,urlcolor=Accent,linkcolor=Accent}
\setlist{leftmargin=1.6em,itemsep=0.3em,topsep=0.35em}
\setcounter{tocdepth}{2}
\renewcommand{\contentsname}{Contents}
\newcommand{\sectionrule}{\par\vspace{-0.5em}\noindent\textcolor{Accent}{\rule{\linewidth}{0.6pt}}\par}
\begin{document}

\begin{center}
{\sffamily\bfseries\color{Accent} TOEFL WORD OF THE DAY\par}
\vspace{8mm}
{\fontsize{42}{48}\selectfont\bfseries prophecy\par}
\vspace{2mm}
{\Large\sffamily\color{Accent} /ˈprɑːfəsi/\par}
\vspace{2mm}
{\small\sffamily Local audio phrase: prophecy\par}
{\small\sffamily Audio file: audios/prophecy.mp3\par}
\vspace{2mm}
{\large\sffamily\color{Muted} countable or uncountable noun\par}
\vspace{5mm}
{\sffamily a prediction believed to have special authority\par}
\vspace{6mm}
{\large A prophecy predicts a future event, traditionally through divine or supernatural knowledge.\par}
\vspace{6mm}
\begin{minipage}{0.84\textwidth}
\itshape “According to the ancient prophecy, a period of peace would follow years of conflict.”
\end{minipage}\par
\vspace{10mm}
\textbf{Date:} September 17th 2026\quad\textbullet\quad \textbf{Level:} B1--mid-B2 TOEFL
\vfill
Created and compiled by \textbf{Amir Kooshky}\par
\vspace{2mm}
\href{https://t.me/KooshkyTOEFL}{Telegram Channel}\quad\textbullet\quad
\href{https://www.instagram.com/kooshkytoefl}{Instagram}\quad\textbullet\quad
\href{https://t.me/KooshkyTOEFL\_pv}{Personal Telegram}
\end{center}
\newpage
\tableofcontents
\newpage

\section{Meanings and usage}
\sectionrule
\subsection{Meaning 1: Prediction with divine or special authority}
\textbf{Part of speech:} countable or uncountable noun

\textbf{IPA:} /ˈprɑːfəsi/

\textbf{Pronunciation phrase:} prophecy

\textbf{Local audio file:} audios/prophecy.mp3

\textbf{Register:} neutral to formal; common in religious, historical, literary, and figurative contexts

\textbf{Definition:} a statement that says what will happen in the future, especially one believed to come from divine or supernatural knowledge

\textbf{In simpler words:} a special or supernatural prediction

\textbf{Typical pattern:} a prophecy about or concerning + event; a prophecy that + clause

\subsubsection{Natural examples}
\begin{enumerate}
  \item According to the ancient prophecy, a period of peace would follow years of conflict.
  \item The story centers on a prophecy that the kingdom will eventually fall.
  \item Many followers believed that the prophecy had been fulfilled.
  \item Historians disagree about when the prophecy was first written down.
  \item The ruler feared a prophecy concerning the end of his reign.
  \item In the novel, an ambiguous prophecy influences nearly every major decision.
  \item The community interpreted the unusual event as the fulfillment of an old prophecy.
  \item A prediction based on scientific evidence is not normally called a prophecy.
\end{enumerate}

\subsubsection{Nuance and implication}
\textbf{Central nuance:} Prophecy often implies divine inspiration, supernatural knowledge, or special authority rather than an ordinary guess about the future.

\textbf{Usual implication:} People may treat the statement as destined to come true, although using the word does not prove that it will.

\textbf{Compared with forecast:} A forecast is normally based on data or observable trends, such as a weather forecast. A prophecy is associated with divine, supernatural, or mysterious knowledge.

\subsubsection{Grammar notes}
\textbf{How to use it:} Use a prophecy for one prediction and prophecies for more than one. Use prophecy without a or an when discussing the activity or tradition generally.

\textbf{What can follow it:} Use about or concerning before a topic, or that before a complete statement.

\textbf{Common preposition:} Common forms are a prophecy about the future, a prophecy concerning the ruler, and a prophecy that the city will fall.

\textbf{Noun use:} It is countable for a particular predicted statement and can be uncountable when referring generally to prophetic speech or tradition.

\subsubsection{Useful patterns}
\textbf{Pattern:} a prophecy about or concerning + noun

\textbf{Use:} Use this to identify the subject of the prophecy.

\textbf{Example:} The manuscript contains a prophecy about the fate of the city.

\textbf{Pattern:} a prophecy that + clause

\textbf{Use:} Use this to state the event that is predicted.

\textbf{Example:} They feared a prophecy that the dynasty would soon end.

\textbf{Pattern:} fulfill or fulfill a prophecy

\textbf{Use:} Use this when later events match what the prophecy predicted.

\textbf{Example:} The character believes that his actions will fulfill the prophecy.

\textbf{Pattern:} a self-fulfilling prophecy

\textbf{Use:} Use this for a belief that helps cause itself to become true.

\textbf{Example:} Expecting the project to fail can become a self-fulfilling prophecy if the team stops making an effort.

\subsubsection{Common wrong usage}
\paragraph{Correction 1}
\textbf{Wrong:} The meteorologist issued a prophecy of heavy rain based on satellite data.

\textbf{Correct:} The meteorologist issued a forecast of heavy rain based on satellite data.

\textbf{Why:} Use forecast for a prediction based on scientific data. Prophecy usually suggests divine or supernatural knowledge.

\paragraph{Correction 2}
\textbf{Wrong:} The priest prophecied that the city would fall.

\textbf{Correct:} The priest prophesied that the city would fall.

\textbf{Why:} The verb is prophesy, and its past form is prophesied.

\paragraph{Correction 3}
\textbf{Wrong:} The prophecy predicted about a coming disaster.

\textbf{Correct:} The prophecy predicted a coming disaster.

\textbf{Why:} Predict takes a direct object, so about is unnecessary here.

\section{Word family}
\sectionrule
\subsection{prophesy}
\textbf{Part of speech:} transitive or intransitive verb

\textbf{IPA:} /ˈprɑːfəsaɪ/

\textbf{Definition:} to state that something will happen in the future, especially through claimed divine or supernatural knowledge

\textbf{Relationship to the headword:} This is the verb corresponding to the noun prophecy.

\textbf{Grammar pattern:} prophesy that + clause, prophesy + event, or prophesy about + topic

\textbf{Usage note:} The final sound is /saɪ/, unlike the /si/ ending of prophecy.

\subsubsection{Examples}
\begin{enumerate}
  \item The text prophesies that the empire will collapse.
  \item Religious leaders prophesied about a coming period of change.
\end{enumerate}

\subsection{prophetic}
\textbf{Part of speech:} adjective

\textbf{IPA:} /prəˈfetɪk/

\textbf{Definition:} relating to prophecy or accurately describing what happens later

\textbf{Relationship to the headword:} This adjective describes a prophecy, a prophet, or a statement that later appears to predict events.

\textbf{Grammar pattern:} prophetic + vision, warning, text, or statement; prove prophetic

\subsubsection{Examples}
\begin{enumerate}
  \item The poem contains several prophetic visions.
  \item Her warning proved prophetic when the system failed a month later.
\end{enumerate}

\subsection{prophet}
\textbf{Part of speech:} countable noun

\textbf{IPA:} /ˈprɑːfɪt/

\textbf{Definition:} a person believed to communicate divine messages or predict the future through divine inspiration

\textbf{Relationship to the headword:} This noun names a person who delivers a prophecy.

\textbf{Usage note:} Do not confuse prophet with profit, although they are pronounced alike in many American accents.

\subsubsection{Examples}
\begin{enumerate}
  \item The tradition regards him as an important prophet.
  \item The prophet warned the people to change their behavior.
\end{enumerate}

\section{Common collocations}
\sectionrule
\subsection{ancient prophecy}
\textbf{Applies to:} Prediction with divine or special authority

\textbf{Meaning:} a prophecy originating in the distant past

\textbf{Example:} The play refers repeatedly to an ancient prophecy about the royal family.

\subsection{fulfill a prophecy}
\textbf{Applies to:} Prediction with divine or special authority

\textbf{Meaning:} cause or appear to cause the predicted event to happen

\textbf{Example:} His unexpected return seemed to fulfill the prophecy.

\subsection{prophecy comes true}
\textbf{Applies to:} Prediction with divine or special authority

\textbf{Meaning:} the event predicted by a prophecy happens

\textbf{Example:} The villagers wait to see whether the prophecy will come true.

\subsection{self-fulfilling prophecy}
\textbf{Applies to:} Prediction with divine or special authority

\textbf{Meaning:} an expectation that influences behavior and helps make the expected result happen

\textbf{Example:} The belief that cooperation was impossible became a self-fulfilling prophecy.

\textbf{Usage note:} This established figurative expression does not require a supernatural meaning.

\section{Synonyms and antonyms}
\sectionrule
\subsection{Close synonyms}
\subsubsection{prediction}
\textbf{Part of speech:} noun

\textbf{Applies to:} Prediction with divine or special authority

\textbf{Shared meaning:} Both state what is expected to happen in the future.

\textbf{Difference:} Prediction is the broad neutral word for any statement about the future. Prophecy usually suggests divine, supernatural, or mysterious authority.

\textbf{Example:} The model's prediction was based on thirty years of climate data.

\subsubsection{forecast}
\textbf{Part of speech:} noun

\textbf{Applies to:} Prediction with divine or special authority

\textbf{Shared meaning:} Both describe a statement about future events.

\textbf{Difference:} A forecast is normally based on data and analysis, especially for weather, economics, or demand. A prophecy is traditionally based on claimed revelation or supernatural insight.

\textbf{Example:} The economic forecast predicts slower growth next year.

\section{Words learners may confuse}
\sectionrule
\subsection{prophesy}
\textbf{Part of speech:} noun versus verb

\textbf{Why learners confuse them:} The spellings differ by only one letter, but the words have different grammatical roles and pronunciations.

\textbf{Key difference:} Prophecy is a noun ending in the sound /si/. Prophesy is a verb ending in /saɪ/.

\textbf{prophecy:} The prophecy describes a future conflict.

\textbf{prophesy:} The speakers prophesy that a conflict will occur.

\subsection{profit}
\textbf{Part of speech:} related family word versus unrelated noun

\textbf{Why learners confuse them:} Prophet and profit are pronounced alike in many American accents.

\textbf{Key difference:} A prophet is a person associated with divine messages; profit is money or another benefit gained from an activity.

\textbf{prophecy:} The prophet delivered a warning.

\textbf{profit:} The company reported a profit this year.

\section{Etymology}
\sectionrule
\textbf{Origin:} Prophecy comes through Old French and Late Latin from Greek prophēteia, referring to the gift or act of interpreting divine will.

\textbf{Development:} The word came to name both divinely inspired speech and a particular statement predicting future events.

\textbf{Memory link:} Connect prophecy, the predicted message, with prophet, the person who delivers it.

\section{Practice exercises}
\sectionrule
\subsection{Word family choice}
Choose the best member of the word family for each blank.

\begin{enumerate}
  \item The ancient text contains a \_\_\_\_\_ about the end of the kingdom.
  \item The warning seemed \_\_\_\_\_ after the predicted shortage occurred.
  \item The speakers \_\_\_\_\_ that a great change will soon occur.
\end{enumerate}

\subsection{Correct or incorrect?}
Decide whether each sentence is correct. Correct the inaccurate ones.

\begin{enumerate}
  \item The legend contains a prophecy that the city will be rebuilt.
  \item The weather prophecy is based on radar observations.
  \item The prophecy eventually came true.
  \item The writer prophecied a period of conflict.
\end{enumerate}

\subsection{Collocation practice}
Complete each sentence with a natural collocation.

\begin{enumerate}
  \item The hero's return appears to fulfill the ancient \_\_\_\_\_.
  \item Their expectation of failure became a self-fulfilling \_\_\_\_\_.
\end{enumerate}

\subsection{Complete answer key}
\subsubsection{Word family choice}
\begin{enumerate}
  \item \textbf{prophecy.} The ancient text contains a prophecy about the end of the kingdom. \emph{Use the noun prophecy after a.}
  \item \textbf{prophetic.} The warning seemed prophetic after the predicted shortage occurred. \emph{Use the adjective prophetic after seemed.}
  \item \textbf{prophesy.} The speakers prophesy that a great change will soon occur. \emph{Use the verb prophesy before the that-clause.}
\end{enumerate}

\subsubsection{Correct or incorrect?}
\begin{enumerate}
  \item \textbf{Correct} \emph{A prophecy that naturally introduces the predicted event.}
  \item \textbf{Incorrect}. The weather forecast is based on radar observations. \emph{Use forecast for a scientific prediction about weather.}
  \item \textbf{Correct} \emph{A prophecy comes true is a natural expression.}
  \item \textbf{Incorrect}. The writer prophesied a period of conflict. \emph{The past form of prophesy is spelled prophesied.}
\end{enumerate}

\subsubsection{Collocation practice}
\begin{enumerate}
  \item \textbf{prophecy.} The hero's return appears to fulfill the ancient prophecy. \emph{Fulfill a prophecy is an established collocation.}
  \item \textbf{prophecy.} Their expectation of failure became a self-fulfilling prophecy. \emph{Self-fulfilling prophecy is an established figurative expression.}
\end{enumerate}

\vfill
\begin{center}
\small\color{Muted} Created and compiled by \textbf{Amir Kooshky}\\
\href{https://t.me/KooshkyTOEFL}{Telegram Channel}\quad\textbullet\quad
\href{https://www.instagram.com/kooshkytoefl}{Instagram}\quad\textbullet\quad
\href{https://t.me/KooshkyTOEFL\_pv}{Personal Telegram}
\end{center}
\end{document}
